\textbf{Phần 1. Trắc nghiệm lựa chọn}
\setcounterpageref{ex}{0}

\begin{ex} %[5H3N15]
Số thích hợp điền vào chỗ chấm để $1\text{ km}^2 = \dots\text{ m}^2$ là
\choice
{$1\,000$}
{\True $1\,000\,000$}
{$10\,000$}
{$100\,000$}
\loigiai{
Ki-lô-mét vuông là diện tích hình vuông cạnh $1\text{ km} = 1\,000\text{ m}$.\\
Do đó $1\text{ km}^2 = 1\,000 \times 1\,000 = 1\,000\,000\text{ m}^2$.\\
\textbf{Đáp án của câu hỏi là B}
}
\end{ex}

\begin{ex} %[5H3N15]
Số thích hợp điền vào chỗ chấm để $1\text{ ha} = \dots\text{ m}^2$ là
\choice
{\True $10\,000$}
{$100$}
{$1\,000$}
{$100\,000$}
\loigiai{
Héc-ta là diện tích hình vuông có cạnh dài $100\text{ m}$.\\
Do đó $1\text{ ha} = 100 \times 100 = 10\,000\text{ m}^2$.\\
\textbf{Đáp án của câu hỏi là A}
}
\end{ex}

\begin{ex} %[5H3N15]
Số thích hợp điền vào chỗ chấm để $1\text{ km}^2 = \dots\text{ ha}$ là
\choice
{$10$}
{$1\,000$}
{\True $100$}
{$10\,000$}
\loigiai{
Ta có $1\text{ km}^2 = 1\,000\,000\text{ m}^2$ và $1\text{ ha} = 10\,000\text{ m}^2$ nên $1\text{ km}^2 = 100\text{ ha}$.\\
\textbf{Đáp án của câu hỏi là C}
}
\end{ex}

\begin{ex} %[5H3H15]
Số thích hợp điền vào dấu hỏi chấm để $2\text{ km}^2 = ?\text{ m}^2$ là
\choice
{$2\,000$}
{$20\,000$}
{$200\,000$}
{\True $2\,000\,000$}
\loigiai{
Vì $1\text{ km}^2 = 1\,000\,000\text{ m}^2$ nên $2\text{ km}^2 = 2\,000\,000\text{ m}^2$.\\
\textbf{Đáp án của câu hỏi là D}
}
\end{ex}

\begin{ex} %[5H3H15]
Số thích hợp điền vào dấu hỏi chấm để $40\,000\text{ ha} = ?\text{ km}^2$ là
\choice
{$4$}
{$40$}
{\True $400$}
{$4\,000$}
\loigiai{
Vì $100\text{ ha} = 1\text{ km}^2$ nên $40\,000\text{ ha} = 40\,000 : 100 = 400\text{ km}^2$.\\
\textbf{Đáp án của câu hỏi là C}
}
\end{ex}

\begin{ex} %[5H3H15]
Đổi $4\text{ km}^2\ 500\text{ m}^2$ ra mét vuông ta được kết quả là
\choice
{$4\,500\text{ m}^2$}
{\True $4\,000\,500\text{ m}^2$}
{$40\,500\text{ m}^2$}
{$400\,500\text{ m}^2$}
\loigiai{
Ta có: $4\text{ km}^2\ 500\text{ m}^2 = 4\,000\,000\text{ m}^2 + 500\text{ m}^2 = 4\,000\,500\text{ m}^2$.\\
\textbf{Đáp án của câu hỏi là B}
}
\end{ex}

\begin{ex} %[5H3H15]
Đổi $5\,000\,000\text{ m}^2$ ra ki-lô-mét vuông ta được kết quả là
\choice
{\True $5\text{ km}^2$}
{$50\text{ km}^2$}
{$0{,}5\text{ km}^2$}
{$500\text{ km}^2$}
\loigiai{
Ta có: $5\,000\,000\text{ m}^2 = 5\,000\,000 : 1\,000\,000 = 5\text{ km}^2$.\\
\textbf{Đáp án của câu hỏi là A}
}
\end{ex}

\begin{ex} %[5H3H15]
Vườn quốc gia Phong Nha -- Kẻ Bàng có diện tích là $123\,326\text{ ha}$. Số đo diện tích vườn quốc gia đó theo đơn vị ki-lô-mét vuông là
\choice
{$12{,}3326\text{ km}^2$}
{$123{,}326\text{ km}^2$}
{\True $1\,233{,}26\text{ km}^2$}
{$12\,332{,}6\text{ km}^2$}
\loigiai{
Ta có: $123\,326\text{ ha} = 123\,326 : 100 = 1\,233{,}26\text{ km}^2$.\\
\textbf{Đáp án của câu hỏi là C}
}
\end{ex}

\begin{ex} %[5H3H15]
Một khu đất hình chữ nhật có kích thước $6\text{ km} \times 4\text{ km}$. Diện tích khu đất đó theo đơn vị héc-ta là
\choice
{$24\text{ ha}$}
{$240\text{ ha}$}
{$24\,000\text{ ha}$}
{\True $2\,400\text{ ha}$}
\loigiai{
Diện tích khu đất là: $6 \times 4 = 24\text{ (km}^2\text{)} = 2\,400\text{ ha}$.\\
\textbf{Đáp án của câu hỏi là D}
}
\end{ex}

\begin{ex} %[5H3V15]
Người ta xây dựng một nhà máy xử lí rác thải trên một khu đất hình vuông có cạnh $200\text{ m}$. Diện tích khu đất đó là bao nhiêu héc-ta?
\choice
{\True $4\text{ ha}$}
{$2\text{ ha}$}
{$40\text{ ha}$}
{$400\text{ ha}$}
\loigiai{
Diện tích khu đất hình vuông là: $200 \times 200 = 40\,000\text{ (m}^2\text{)}$.\\
Đổi $40\,000\text{ m}^2 = 4\text{ ha}$.\\
\textbf{Đáp án của câu hỏi là A}
}
\end{ex}

\begin{ex} %[5H3V15]
Các kĩ sư dự kiến lắp đặt nhà máy điện mặt trời trên một trong ba mảnh đất: Mảnh A hình chữ nhật $7\text{ km} \times 3\text{ km}$; Mảnh B hình vuông cạnh $5\text{ km}$; Mảnh C hình chữ nhật $6\text{ km} \times 4\text{ km}$. Kĩ sư sẽ chọn mảnh đất có diện tích lớn nhất là
\choice
{Mảnh A}
{\True Mảnh B}
{Mảnh C}
{Cả ba mảnh đất có diện tích bằng nhau}
\loigiai{
Diện tích mảnh A: $7 \times 3 = 21\text{ km}^2$.\\
Diện tích mảnh B: $5 \times 5 = 25\text{ km}^2$.\\
Diện tích mảnh C: $6 \times 4 = 24\text{ km}^2$.\\
Vì $25\text{ km}^2 > 24\text{ km}^2 > 21\text{ km}^2$ nên chọn Mảnh B.\\
\textbf{Đáp án của câu hỏi là B}
}
\end{ex}

\begin{ex} %[5H3V15]
Nối số đo thích hợp với mỗi địa danh: Tỉnh Thừa Thiên Huế, Vườn quốc gia Bạch Mã, Điện Thái Hoà. Số đo diện tích phù hợp với Vườn quốc gia Bạch Mã là
\choice
{$5\,054\text{ km}^2$}
{$1\,360\text{ m}^2$}
{$100\text{ ha}$}
{\True $22\,031\text{ ha}$}
\loigiai{
\begin{itemize}
    \item Điện Thái Hoà là một công trình kiến trúc $\rightarrow 1\,360\text{ m}^2$.
    \item Vườn quốc gia Bạch Mã là khu bảo tồn thiên nhiên lớn $\rightarrow 22\,031\text{ ha}$ (tương đương khoảng $220\text{ km}^2$).
    \item Tỉnh Thừa Thiên Huế là đơn vị hành chính cấp tỉnh $\rightarrow 5\,054\text{ km}^2$.
\end{itemize}
\textbf{Đáp án của câu hỏi là D}
}
\end{ex}

\textbf{Phần 2. Trắc nghiệm đúng sai}
\setcounterpageref{ex}{0}

\begin{ex} %[5H3H15]
Xét tính đúng sai của các phép đổi đơn vị đo diện tích sau:
\choiceTF
{\True $3\text{ km}^2 = 300\text{ ha}$.}
{$64\,800\text{ ha} = 64{,}8\text{ km}^2$.}
{\True $6\,450\,000\text{ m}^2 = 6{,}45\text{ km}^2$.}
{\True $2\text{ ha } 500\text{ m}^2 = 20\,500\text{ m}^2$.}
\loigiai{
\begin{itemize}
    \item $3\text{ km}^2 = 3 \times 100 = 300\text{ ha}$.
    \item $64\,800\text{ ha} = 64\,800 : 100 = 648\text{ km}^2 \neq 64{,}8\text{ km}^2$.
    \item $6\,450\,000\text{ m}^2 = 6\,450\,000 : 1\,000\,000 = 6{,}45\text{ km}^2$.
    \item $2\text{ ha } 500\text{ m}^2 = 20\,000 + 500 = 20\,500\text{ m}^2$.
\end{itemize}
\begin{itemize}
    \item a) \textbf{Đúng}.
    \item b) \textbf{Sai}: $64\,800\text{ ha} = 648\text{ km}^2$.
    \item c) \textbf{Đúng}.
    \item d) \textbf{Đúng}.
\end{itemize}
}
\end{ex}

\begin{ex} %[5H3V15]
Một khu bảo tồn thiên nhiên có diện tích là $7\,500\text{ ha}$. Xét tính đúng sai của các khẳng định sau:
\choiceTF
{\True Diện tích khu bảo tồn đổi ra ki-lô-mét vuông là $75\text{ km}^2$.}
{Diện tích khu bảo tồn đổi ra mét vuông là $750\,000\text{ m}^2$.}
{\True Nếu mở rộng thêm $25\text{ km}^2$ thì diện tích mới là $100\text{ km}^2$ (hay $10\,000\text{ ha}$).}
{Diện tích khu bảo tồn này lớn hơn diện tích một tỉnh rộng $5\,000\text{ km}^2$.}
\loigiai{
\begin{itemize}
    \item $7\,500\text{ ha} = 7\,500 : 100 = 75\text{ km}^2$.
    \item $7\,500\text{ ha} = 7\,500 \times 10\,000 = 75\,000\,000\text{ m}^2 \neq 750\,000\text{ m}^2$.
    \item Diện tích mới: $75 + 25 = 100\text{ km}^2 = 10\,000\text{ ha}$.
    \item $75\text{ km}^2 < 5\,000\text{ km}^2$ nên diện tích khu bảo tồn nhỏ hơn diện tích tỉnh đó.
\end{itemize}
\begin{itemize}
    \item a) \textbf{Đúng}.
    \item b) \textbf{Sai}: Kết quả là $75\,000\,000\text{ m}^2$.
    \item c) \textbf{Đúng}.
    \item d) \textbf{Sai}: $75\text{ km}^2 < 5\,000\text{ km}^2$.
\end{itemize}
}
\end{ex}

\textbf{Phần 3. Trả lời ngắn}
\setcounterpageref{ex}{0}

\begin{ex} %[5H3H15]
Đổi $64\,800\text{ ha}$ ra đơn vị ki-lô-mét vuông.
\shortans{648}
\loigiai{
\textbf{Đáp số:} $648$\\
$64\,800\text{ ha} = 64\,800 : 100 = 648\text{ km}^2$.
}
\end{ex}

\begin{ex} %[5H3H15]
Đổi $3\text{ km}^2$ ra héc-ta.
\shortans{300}
\loigiai{
\textbf{Đáp số:} $300$\\
$3\text{ km}^2 = 3 \times 100 = 300\text{ ha}$.
}
\end{ex}

\begin{ex} %[5H3H15]
Một khu rừng có diện tích $6\,450\,000\text{ m}^2$. Hỏi diện tích khu rừng đó bằng bao nhiêu ki-lô-mét vuông? (Nhập kết quả dưới dạng số thập phân).
\shortans{6,45}
\loigiai{
\textbf{Đáp số:} $6{,}45$\\
$6\,450\,000\text{ m}^2 = 6\,450\,000 : 1\,000\,000 = 6{,}45\text{ km}^2$.
}
\end{ex}

\begin{ex} %[5H3V15]
Khu đất hình vuông có cạnh $200\text{ m}$. Hỏi diện tích khu đất đó là bao nhiêu héc-ta?
\shortans{4}
\loigiai{
\textbf{Đáp số:} $4$\\
Diện tích là $200 \times 200 = 40\,000\text{ m}^2 = 4\text{ ha}$.
}
\end{ex}

\begin{ex} %[5H3V15]
Mảnh đất hình chữ nhật có chiều dài $5\text{ km}$ và chiều rộng $3\text{ km}$. Hỏi diện tích mảnh đất đó bằng bao nhiêu héc-ta?
\shortans{1500}
\loigiai{
\textbf{Đáp số:} $1500$\\
Diện tích mảnh đất là: $5 \times 3 = 15\text{ (km}^2\text{)} = 1\,500\text{ ha}$.
}
\end{ex}

\begin{ex} %[5H3V15]
Vườn quốc gia Tràm Chim có diện tích $7\,500\text{ ha}$, Vườn quốc gia Bạch Mã có diện tích $22\,031\text{ ha}$. Hỏi diện tích Vườn quốc gia Bạch Mã lớn hơn diện tích Vườn quốc gia Tràm Chim bao nhiêu héc-ta?
\shortans{14531}
\loigiai{
\textbf{Đáp số:} $14531$\\
Hiệu diện tích là: $22\,031 - 7\,500 = 14\,531\text{ (ha)}$.
}
\end{ex}
